\chapter*{Введение}                         % Заголовок
\addcontentsline{toc}{chapter}{Введение}    % Добавляем его в оглавление

\newcommand{\actuality}{}
\newcommand{\progress}{}
\newcommand{\aim}{{\textit\aimTXT}}
\newcommand{\tasks}{\textit{\tasksTXT}}
\newcommand{\theoretical}{\textit{\theoreticalTXT}} % Added by ASM
\newcommand{\novelty}{\textit{\noveltyTXT}}
\newcommand{\influence}{\textit{\influenceTXT}}
\newcommand{\methods}{\textit{\methodsTXT}}
\newcommand{\defpositions}{\textit{\defpositionsTXT}}
\newcommand{\pasport}{{\textit\pasportTXT}}
\newcommand{\reliability}{\textit{\reliabilityTXT}}
\newcommand{\probation}{\textit{\probationTXT}}
\newcommand{\reliaprobation}{\textit{\reliaprobationTXT}}
\newcommand{\contribution}{\textit{\contributionTXT}}
\newcommand{\publications}{\textit{\publicationsTXT}}

\input{common/characteristic} % Характеристика работы по структуре во введении и в автореферате не отличается (ГОСТ Р 7.0.11, пункты 5.3.1 и 9.2.1), потому её загружаем из одного и того же внешнего файла, предварительно задав форму выделения некоторым параметрам

\textbf{Объем и структура работы.} Диссертация состоит из~введения,
\formbytotal{totalchapter}{глав}{ы}{}{}, \label{totalchapter}
заключения и
\formbytotal{totalappendix}{приложен}{ия}{ий}{ий}. 
%% на случай ошибок оставляю исходный кусок на месте, закомментированным
%Полный объём диссертации составляет  \ref*{TotPages}~страницу
%с~\totalfigures{}~рисунками и~\totaltables{}~таблицами. Список литературы
%содержит \total{citenum}~наименований.
%
Полный объём диссертации составляет
\formbytotal{TotPages}{страниц}{у}{ы}{}, включая
\formbytotal{totalcount@figure}{рисун}{ок}{ка}{ков} и %\lable{figsss}
\formbytotal{totalcount@table}{таблиц}{у}{ы}{}.
Список литературы содержит
\formbytotal{citenum}{наименован}{ие}{ия}{ий}.
